\chapter{Conclusion and Future Work of KD fermions}\label{chapter:FutureWork_KD}

\chapterabstract{In this chapter, we will conclude this part of the thesis, and then outline several compelling directions for future research regarding Kähler-Dirac (KD) fermions. First, we discuss their application to lattice field theory, demonstrating how the two-sheeted spacetime structure and the inclusion of four KD fermion copies facilitate the construction of an anomaly-free chiral gauge theory. Next, we explore the profound geometric connections between KD fields and twistor theory. Specifically, we detail how the internal symmetries and two-sheeted topology of the KD framework map naturally onto the conformal group and the self-dual/anti-self-dual (SD/ASD) decomposition within twistor space. Unifying these two mathematical frameworks may provide a natural resolution to long-standing puzzles, such as the existence of exactly three generations of fermions. Finally, we examine the potential integration of KD fermions into superstring theory. The inherent two-sheeted structure might theoretically enable the cancellation of ghosts, paving the way for a consistent 4-dimensional string theory. Because knots and braids are topologically non-trivial strictly in three spatial dimensions, this approach establishes a compelling link to Kauffman's topological Preon model. In this framework, Standard Model (SM) particles are represented by distinct braids of strings, offering a novel geometric paradigm for resolving many of the remaining mysteries within the SM. }

\section{Conclusion}

This part began by introducing Kähler-Dirac (KD) fermions, detailing their applications in lattice gauge theory, the mathematical mapping between their differential form and spinor formulations, and their physical properties in Lorentzian signature (\cref{chapter:Intro_Basics}). We demonstrated that the additional $\gamma^0$ required to construct a Lorentz-invariant Lagrangian inevitably causes half of the field components to exhibit negative energies and negative norms. To resolve this physical pathology, a two-sheeted, $PT$-mirrored spacetime geometry was introduced, wherein the negative-energy states naturally reside on the $PT$-reversed sheet. To prevent unphysical mixing between the two sheets—which would explicitly violate causality—we exploited the internal symmetry degrees of freedom of the KD field. This rigorous restriction not only forces the KD fields to transform strictly as fermions (Grassmann-valued fields), but also constrains the internal symmetry to the compact subgroup $U(2)\times U(2)$. This symmetry dynamically manifests as the electroweak force acting in a mirror-symmetric fashion across the two sheets. To maintain consistency under charge conjugation, a KD-Majorana condition is strictly required. This condition dictates that particles on the second sheet are exactly the antiparticles of their first-sheet counterparts, mathematically ensuring that the two sheets remain perfectly $CPT$-symmetric, even at the quantum scale (\cref{chapter:two-sheet_KD_Majorana}).

In \cref{chapter:SM_Cosmology}, we explored the profound implications of this framework for both the Standard Model (SM) of particle physics and cosmology. All 16 Weyl fermions of a single generation are elegantly accommodated within four copies of the KD field (representing one lepton and three quark colors). The complete SM gauge group, $SU(3)\times SU(2)\times U(1)$, naturally emerges from the unitary transformations among the three quark copies combined with the KD internal symmetry $U(2)\times U(2)$. Furthermore, the fermionic kinetic terms, Yukawa couplings, and right-handed Majorana mass terms were formulated explicitly within this unified KD framework. We also discussed the cosmological implications of this model, highlighting its seamless integration with the $CPT$-symmetric universe paradigm and its direct application to the black hole mirror hypothesis.

Finally, in \cref{chapter:Lag_Euclidean}, we examined the Lagrangian of KD fermions in Euclidean signature. Directly Wick rotating the Lorentzian Lagrangian into its Euclidean counterpart via von Nieuwenhuizen's continuous transformation method allowed us to construct a theory that is simultaneously Hermitian and chiral. Inspired by this emergent Euclidean formulation, we explored the theoretical viability of alternative Lorentz-invariant Lagrangians in Minkowski space. Our rigorous analysis demonstrates that, despite explicitly different mathematical formulations, the underlying physical paradigms—including the necessity of the two-sheeted structure and the KD-Majorana constraints—remain fundamentally invariant.

% - Summarize the finding.

% % Editor's note: Capitalized the subsection title for consistency with thesis formatting.
% \subsection{Theoretical Explanation for the Absence of Mixing Terms}
% % - Future work 1 -- theoretical explanation why there is no mixing terms: in the thesis, we show that buy choosing internal symmetry V, we can cancel $U_\Lambda^{-1}$ and spin connection term acting from the right, and go back to conventional Dirac spinors (it is important not just for matching fermions spin-1/2 property, but also prevent mixing between forward and backward sheets). However, do we have a fundamental reason why they couldn't exist? To answer this questions, maybe we can resort to fermions interaction in curved spacetime. Kostas' work shows that IC of gravity should be Neumann BC. So if fermions interact with gravity, it would also set constraint on IC of fermions (prevent mixing transformation), which might propagate to later universe.

% Throughout this thesis, we have asserted that fermions residing on the two spacetime sheets must not mix; otherwise, causality would be violated, in stark conflict with observation. However, imposing this non-mixing constraint phenomenologically appears somewhat ad hoc. We must therefore ask: is there a fundamental dynamical reason that prevents this mixing from occurring? 

% A promising potential answer lies in the interaction between fermions and gravity. Kostas Tzanavaris' recent work \cite{Tzanavaris:2026ujl} demonstrates that the initial conditions of gravity at the initial singularity must strictly satisfy Neumann boundary conditions. If fermions are dynamically coupled to gravity, this gravitational boundary condition could impose strict constraints on the fermionic initial conditions as well, effectively forbidding mixing transformations. These constraints might then propagate forward into the late universe. For example, while spatial rotations preserve the isolation of the sheets, a Lorentz boost would theoretically mix them. Yet, under a Lorentz boost, Neumann boundary conditions (which require specific time derivatives to vanish) are no longer satisfied. Similar rigid constraints might apply to the spin connection in highly curved spacetime. We intend to explore this geometric decoupling mechanism rigorously in future work.

\section{Future work}

In this section, I will discuss three direction of future work, including its application in chiral gauge theory, twistor theory, and 4D supersting theory.

% Editor's note: Capitalized the subsection title for consistency.
\subsection{Applications to Lattice Chiral Gauge Theory}
% - Future work2 -- Lattice field theory: in the thesis, we show a Euclidean action which is both Hermitian and chiral. Moreover, we get the same conclusion as Barret work: two sheeted is needed to solve problems in lattice field theory. Hope can construct a problem free mathematical framework for simulating fermions on lattice.

As mentioned in the introduction \cref{sec:KD_intro}, KD fermions has profound applications in lattice gauge theory. Catterall's recent papers \cite{Catterall:2022jky, Catterall:2023nww} state that to construct an anomaly-free Kähler-Dirac (KD) chiral gauge theory on a curved spacetime lattice, one requires mirror fermions (which maps perfectly onto our two-sheeted topological picture). These mirror fermions are necessary to remove the phase ambiguities arising from an unequal number of chiral fermions in general curved geometries—a phenomenon mathematically analogous to the usual $U(1)$ chiral anomaly of a single Weyl field. Moreover, exactly $0\text{ mod }4=4,8,12,\dots$ copies of KD fermions paired with mirror KD fermions are required to cancel the anomaly originating from the global $U(1)$ symmetry, making 4 copies the minimal viable theory. This comfortably accommodates all 16 Weyl spinors of a single Standard Model generation. 

In other words, the two-sheeted spacetime picture featuring mirror (KD-Majorana) fermions across 4 copies is not merely an ad hoc solution for removing negative-norm states and explaining the SM; it is a fundamental mathematical requirement for constructing an anomaly-free chiral gauge theory. In Catterall's work \cite{Catterall:2023nww}, explaining the observational absence of these mirror fermions required invoking an external mechanism (such as coupling the fermions to a scalar field) to artificially gap them out. In this thesis, we point out that these mirror fermions are physically nothing more than the standard antiparticles residing on the backward causal sheet. This resolves the observational issue elegantly, requiring no artificial gapping or spontaneous symmetry breaking.

Furthermore, an upcoming companion paper by Latham Boyle and Vatsalya Vaibhav \cite{BoyleVaibhav} provides the rigorous mathematical details of how to quantize the KD field directly on a lattice. To tackle the longstanding problem of formulating a chiral Lagrangian in Euclidean signature, they build upon Barrett's recent work \cite{Barrett2024}. They demonstrate that by doubling the Hilbert space—i.e., introducing KD fields that live on two sheets of a 4D spacetime—one can construct a strictly chiral Lagrangian of the form $\big<J\bar{\Psi}_L,K\Psi_L\big>+\big<J\bar{\Psi}_R,K\Psi_R\big>$. Interestingly, this lattice construction perfectly mirrors the two-sheeted topological picture we independently derived in the continuum limit (\cref{eq:Lagrangian_KD_Euclidean}). Ultimately, the two-sheeted spacetime is not merely a trick for excising negative-norm states; it is a rigorous mathematical prerequisite for defining a valid chiral Lagrangian for chiral gauge theories in Euclidean signature.

% Editor's note: Capitalized the subsection title for consistency.
\subsection{Geometric Connections to Twistor Theory}
% - Future work3 -- KD fermions's geometry might link to twistor theory. Ex: they both have U(2,2), two sheeted of spcetime <-> twistor space and dual twistor space. --> combination of them might solve the problems such as why three generations of fermions.

The unique symmetry and geometric structure of the Lorentzian KD action strongly suggests a natural geometric connection to twistor space \cite{Penrose:1985bww, Penrose:1986ca, adamo2018lecturestwistortheory}, whose fundamental conformal symmetry group is $SU(2,2)$ (aligned with the KD internal symmetry group $U(2,2)$). The intrinsic division into self-dual and anti-self-dual spaces and fields may be deeply related to our two-sheeted spacetime model. It is particularly worth highlighting the proposed connection between twistor space and the Standard Model outlined in \cite{woit2021euclideantwistorunification}, which poses the open question of whether situating KD fermions within twistor theory might naturally explain the existence of exactly three generations of fermions. Here are the outlines of overlaps between KD fermions and twistor theory, which are interesting to explore:

\begin{enumerate}
    \item \textbf{The "Two Sheets" as SD/ASD Decomposition:} In twistor theory, physical fields on complexified spacetime naturally decompose into Self-Dual (SD) and Anti-Self-Dual (ASD) sectors. Each sector can evolve independently—a cornerstone of twistor constructions for both gauge and gravitational theories. We hypothesize that our "two sheets" picture is compatible with this fundamental SD/ASD splitting. 
    \item \textbf{Internal Symmetry as the Conformal Group:} The KD algebra exhibits an internal $U(2, 2)$ symmetry. It seems to be a remarkable physical coincidence that the fundamental symmetry group of twistor space is precisely $SU(2, 2)$ (the conformal group, acting as a double cover of $SO(4, 2)$). Could the "internal" symmetries of the KD field simply be the action of the conformal group on the field components, made manifest when lifted to the twistor picture? This would elegantly unify spacetime and internal symmetries into a single geometric origin.
    % Editor's note: Fixed the raw unicode bullet point to the standard LaTeX \bullet command inside math mode.
    \item \textbf{The Complex Structure of Twistor Space:} Twistor space is a complex spacetime, which unifies Euclidean and Lorentzian signatures as different real slices of the same overarching complex manifold, potentially resolving the lingering tension between Euclidean lattice methods and real Lorentzian physics. If the rsult consistent with Wick rotation discussed in \cref{sec:extend_KD}, it would be a strong geometrical support of the KD Euclidean theory.
\end{enumerate}

This unified perspective raises several intriguing, albeit highly speculative, questions:
\begin{enumerate}
    \item \textbf{Explainin the SM:} Peter Woit \cite{woit2021euclideantwistorunification} pointed out that the SM gauge group and the specific fermion content of a single generation are natively embedded in the geometry of twistor space. Because these properties are also embedded in the geometry of KD fields, could these two pictures be perfectly unified when lifted into twistor space? Furthermore, as suggested by Woit, could this unified geometric framework finally explain why exactly three generations of fermions exist?
    \item \textbf{Discrete Twistor Spaces:} Rather than discretizing spacetime directly, would it be mathematically more natural to define a KD-like operator on a discrete twistor space? For instance, \citet{ShapiroGeorge2013} demonstrated how to construct a twistor lattice. Could structures like "discrete conjugate nets" provide a lattice that preserves a discrete subgroup of the conformal group while naturally supporting a complex structure? If such a complex lattice gauge theory is successfully built, one can simulate fermions on a lattice and then map to the Euclidean or Lorentzian signature simply by taking real slices.
\end{enumerate}

However, significant mathematical challenges remain to be tackled. For example, the clean separation of SD and ASD sectors breaks down under full General Relativity (which fundamentally mixes SD and ASD components). In other words, $\bar\partial^2\neq0$, causing unwanted mixing terms in dual $\mathbb{PT}$ space. Similarly, for the Kähler-Dirac field, consider $D_\mu=\partial_\mu+[\omega_\mu,\cdot]$ the right acting spin connection in strongly curved spacetime actively mixes the two sheets. This points to the compelling possibility that KD fields satisfying the KD-Majorana condition actually reside in Ambitwistor space—the geometric intersection between twistor space ($Z$) and dual twistor space ($W$), defined by the incidence relation $Z\cdot W=0$. 

While we have demonstrated that one can carefully choose the internal symmetry of a KD field such that curved spacetime does not mix the two sheets \cref{subsec:solution_internal_symmetry}, it is not yet clear whether a similar geometric decoupling can be naturally achieved within Ambitwistor space. This warrants intensive further investigation.

% Editor's note: Capitalized the subsection title for consistency.
\subsection{Towards a Four-Dimensional KD Superstring Theory}

% - Future work4 -- construct a new superstring theory via KD fermions. current superstring theory have 10 dimensions. The question is whether we can construct a theory by KD fermions, which is ghost free in just two sheeted of 4D spacetime? And since there are only 3d space, tangle and knots are distinguishable. By applying the interesting Preons braid theory by Louis Kauffman, might open to a new door for understanding fundamental physics.

String theory is a prominent candidate for a unified theory of quantum gravity. By replacing fundamental point particles with 1D strings—such that 1D worldlines become 2D worldsheets—Feynman diagrams no longer contain singular interaction vertices. Instead, vertices are smoothed out into finite intersections of worldsheets, transforming point-like interactions into the topology of manifolds. Consequently, one no longer needs to mathematically cancel non-physical infinities (via renormalization) when integrating over singular intersection paths.

However, a well-known theoretical constraint arises: to ensure the 2D worldsheet metric is conformally invariant (anomaly-free), the target spacetime must possess 26 dimensions for bosonic strings, or 10 dimensions for superstrings. This not only conflicts starkly with macroscopic observation (which reveals only a 4-dimensional spacetime) but also introduces immense theoretical complexity. While string theorists argue that this dimensionality problem is resolved by compactification—with the additional dimensions curling up into tightly bound, unobservable microscopic manifolds (e.g., Calabi-Yau manifolds)—and that these extra degrees of freedom are actually a useful feature for explaining dark matter, dark energy, inflation, and Grand Unification, the sheer landscape of solutions remains problematic.

Applying Occam's razor, we must ask: are there simpler solutions to these fundamental questions? If we rigorously analyze the mathematical origin of these additional dimensions, we find three hidden assumptions in the standard derivation:
\begin{enumerate}
    \item \textbf{A single fundamental string type:} This assumption is easily relaxed by introducing multiple distinct types of strings (e.g., carrying internal "color" indices like red, blue, and green).
    \item \textbf{A standard Riemannian target space:} Barrett's recent work \cite{Barrett2024} demonstrates that for a strictly chiral fermionic theory, the Euclidean spacetime dimension must be $2 \text{ mod } 8$ (e.g., 2, 10, 18, 26). Standard string theory targets 10 or 26. However, Barrett also proved that a chiral theory can exist in 4D if the Hilbert space is doubled—precisely mapping to our two-sheeted 4D spacetime.
    \item \textbf{Scale (conformal) invariance:} This symmetry must be satisfied on the 2D worldsheet to decouple unphysical degrees of freedom; we retain this assumption as fundamental.
\end{enumerate}

The central question is therefore: could we construct a ghost-free string theory situated purely in a two-sheeted 4D spacetime by incorporating KD fermions on the worldsheet? 

The idea of constructing a KD-based string theory was first explored in by Jourjine \cite{Jourjine1986}, who pointed out that the resulting Lagrangian identically cancels to zero. A subsequent comment \cite{Bullinaria1987CommentOJ} demonstrated that this cancellation is directly caused by the additional $\gamma^0$ in the KD Lagrangian, which forces half of the field components to acquire a negative sign. This is essentially identical to the phenomena evaluated throughout this thesis. In his paper, Bullinaria discussed the possibility of decoupling or excising these negative states (which he viewed as non-physical ghosts) from the Lagrangian, finding the mathematical task highly challenging. However, our work demonstrates that this cancellation physically signifies the necessity of a two-sheeted, $PT$-reversed spacetime. By adopting this topological picture, it may be possible to move past Bullinaria's roadblocks and successfully construct a consistent KD string theory. 

The path forward, however, requires resolving several theoretical hurdles. Conventional superstring theory relies on supersymmetry between its fermionic (2D Dirac spinor) and bosonic (spacetime coordinate) components. Because the KD framework inherently contains more spinors (doubling the fermionic degrees of freedom in 2D), does this demand a correspondingly enlarged spacetime symmetry? Do the worldsheet fields map target coordinates onto two distinct sheets of spacetime? These questions must be formally addressed to construct a viable KD-string Lagrangian, from which one can then rigorously evaluate ghost cancellation.

Furthermore, if this 4D framework proves successful, it unlocks a profound topological feature: because the spatial manifold is strictly 3-dimensional, strings can form stable, distinguishable knots and braids. In standard string theories (where the spatial dimension is $\geq 4$), knots are trivial because strings can always be untangled via continuous transformations. In a 4D spacetime, however, topology matters. For example, a worldsheet could be formed by a trefoil knot propagating along the time axis, or by three distinct strings (e.g., red, green, and blue) braiding around one another. 

A series of papers by Louis H. Kauffman regarding the topological Preon model \cite{bilsonthompson2006topologicalmodelcompositepreons, Chester2025} suggests that all SM particles—including all gauge bosons and fermions—might be physically constructed by three "colored" ribbons braided together. This braiding topology can successfully explain all known fermions with their correct chiralities and hypercharges, reproduce the $SU(3)\times U(1)_Y$ gauge group, and dynamically predict composite particle interactions 
% Editor's note: Fixed the raw unicode arrows and greek letters in the reaction equation to proper LaTeX math commands.
(e.g., $n \to W_- + p^+ \to e_L + \bar{\nu}_R + p^+$). 

Synthesizing this topological Preon model with a rigorous KD string theory may yield a remarkably powerful and predictive framework. For example, one of the main challenges of the Preon model is its inability to geometrically explain the weak force $SU(2)$ or why it couples exclusively to left-handed particles. As demonstrated in this thesis, this chiral gage behavior emerges naturally from the internal symmetry of the KD field, which might explain the puzzle.


%%%%%%%%%%%%%%%%
% Other applications
%%%%%%%%%%%%%%%%

% \subsection{Other Applications: Fermion-Rotor Systems and Monopole Scattering}

% The two-sheeted topological picture may find vital applications in other domains of theoretical physics. Consider, for example, the long-standing puzzle of monopole scattering: when a chiral fermion scatters off a magnetic monopole, it theoretically bounces back with radically altered quantum numbers (e.g., flipped chirality or fractional charge). To resolve this paradox, \citet{Loladze:2025jsq} considered a "fermion-rotor" toy model, demonstrating that a quantum rotor acts as a "twist operator" that dynamically alters the quantum numbers of the fermions, effectively dressing them in a topological string to preserve both symmetries and overall unitarity. 

% Alternatively, the polyform, two-sheeted geometry proposed in this thesis might naturally resolve the monopole scattering paradox without needing to artificially couple the system to an external "rotor." Just as the KD-Majorana condition inherently pairs a left-handed particle on the forward sheet with a right-handed antiparticle on the backward sheet, a fermion striking a fundamental monopole might be mathematically understood as traversing this $CPT$-mirror boundary. This topological traversal would naturally and deterministically explain why the outgoing particle emerges with flipped or altered quantum numbers.

% It is important to emphasize, however, that the $CPT$-symmetric universe model does not natively predict the existence of macroscopic magnetic monopoles. Magnetic monopoles are typically produced when a Grand Unified Theory (GUT) spontaneously breaks down to a smaller symmetry group containing electromagnetism (a $U(1)$ factor). Because our framework relies solely on Pati-Salam unification—whose specific symmetry-breaking pathways do not topologically necessitate monopoles—they are absent from the baseline cosmology. The monopole scattering problem is mentioned here merely to illustrate how the two-sheeted topological paradigm provides a powerful conceptual tool capable of resolving complex boundary phenomena across various fields of physics.

\printbibliography[heading=subbibintoc, title={References for Chapter \thechapter}]